\section[Cote 41 : continuité de la section générique et ouverture (proposition 1)]{Cote 41 : continuité de la section générique et ouverture (proposition 1)\protect\verifie{Folder41}}\label{sec:41}

La cote 41, \emph{Construction de morphismes étales (Chap. V) : notes manuscrites (s.d.)}, datée par l'inventaire \q{[années 1960-1970]}, est une seule section de rédaction en treize pages. Elle pose, page~3, la notion de section \emph{monogène} --- une section qui choisit continûment le point générique d'une composante géométriquement irréductible de chaque fibre --- et la question de la représentabilité du foncteur qu'elles forment ; elle construit ensuite, sous des conditions de spécialisation, un revêtement étale séparé qui la représente (pages~4 à~10). Page~11, elle revient au cas où les fibres elles-mêmes sont irréductibles, et y écrit une \q{Prop.~1°} en trois conditions, à laquelle s'ajoute, encadrée, une condition \q{(iii bis)} ; l'hypothèse de finitude est portée sur la condition~(ii), \q{$f$ soit ouvert [\dots] si $Y$ est loc.\ noeth., $f$ de type fini}, et la seule formule lisible de la démonstration est $\varphi^{-1}(U) = f(U)$ (\lot{41}{1}{11}).

\begin{enonce}[Proposition 1]
Soit $f : X \to Y$ un morphisme de type fini dont les fibres sont irréductibles et non vides, et soit $\varphi(y)$ le point générique de $X_y$. On considère les conditions :
\begin{enumerate}
\item[(i)] $\varphi : Y \to X$ est continue ;
\item[(ii)] $f$ est ouvert ;
\item[(iii)] $y \in \overline{\{y'\}} \Rightarrow \varphi(y) \in \overline{\{\varphi(y')\}}$ ;
\item[(iii bis)] pour tous $y \in \overline{\{y'\}}$ et $x$ au-dessus de $y$, il existe $x'$ au-dessus de $y'$ tel que $x \in \overline{\{x'\}}$.
\end{enumerate}
Alors (i)~$\Leftrightarrow$~(ii) sans autre hypothèse, et (ii)~$\Rightarrow$~(iii)~$\Leftrightarrow$~(iii bis) ; si $f$ est localement de présentation finie --- par exemple $Y$ localement noethérien ---, les quatre conditions sont équivalentes.
\end{enonce}

On précise ce que veut dire, sur les espaces sous-jacents, que $\varphi(y)$ est le point générique de la fibre.

\begin{definition}
Soient $X$, $Y$ des espaces topologiques et $f : X \to Y$ une application quelconque. On dit que $\varphi : Y \to X$ est une \emph{section générique} de $f$ si, pour tout $y \in Y$, on a $f(\varphi(y)) = y$ et $x \in \overline{\{\varphi(y)\}}$ pour tout $x \in f^{-1}(y)$.
\end{definition}

\noindent Cela revient à dire que $\varphi(y)$ est un point générique de la fibre $f^{-1}(y)$ munie de la topologie induite, puisque l'adhérence de $\{\varphi(y)\}$ dans $f^{-1}(y)$ est la trace sur $f^{-1}(y)$ de son adhérence dans $X$ ; une fibre qui a un point générique est irréductible et non vide. Pour un morphisme de schémas, la fibre schématique $X_y$ est homéomorphe au sous-espace $f^{-1}(y)$ ; ce fait classique n'est pas démontré ici. Dans les lemmes qui suivent, $f$ est une application quelconque entre espaces topologiques, qu'on ne suppose même pas continue, et $\varphi$ une section générique de $f$.

\begin{lemme}\label{lem:41-formule}
Pour tout ouvert $U$ de $X$, $\varphi^{-1}(U) = f(U)$.
\end{lemme}

\begin{proof}
Si $\varphi(y) \in U$, alors $y = f(\varphi(y)) \in f(U)$. Inversement, soient $x \in U$ et $y = f(x)$ ; alors $x \in \overline{\{\varphi(y)\}}$, et l'ouvert $U$, qui contient un point de l'adhérence de $\{\varphi(y)\}$, contient $\varphi(y)$ : $y \in \varphi^{-1}(U)$.
\end{proof}

\begin{lemme}\label{lem:41-top}
On a (i)~$\Leftrightarrow$~(ii), (i)~$\Rightarrow$~(iii), et (iii)~$\Leftrightarrow$~(iii bis).
\end{lemme}

\begin{proof}
(i)~$\Leftrightarrow$~(ii). Par le lemme~\ref{lem:41-formule}, pour tout ouvert $U$ de $X$, l'ensemble $\varphi^{-1}(U)$ est ouvert si et seulement si $f(U)$ l'est ; or $\varphi$ est continue si et seulement si le premier l'est pour tout $U$, et $f$ est ouverte si et seulement si le second l'est pour tout $U$.

(i)~$\Rightarrow$~(iii). Une application continue $\varphi$ envoie l'adhérence de $\{y'\}$ dans l'adhérence de $\{\varphi(y')\}$.

(iii)~$\Rightarrow$~(iii bis). Soient $y \in \overline{\{y'\}}$ et $x$ au-dessus de $y$. Alors $x \in \overline{\{\varphi(y)\}}$ par définition de $\varphi$, et $\varphi(y) \in \overline{\{\varphi(y')\}}$ par (iii), donc $x \in \overline{\{\varphi(y')\}}$ ; et $x' = \varphi(y')$ est au-dessus de $y'$.

(iii bis)~$\Rightarrow$~(iii). Soit $y \in \overline{\{y'\}}$. Appliquons (iii bis) à $x = \varphi(y)$, qui est au-dessus de $y$ : il existe $x'$ au-dessus de $y'$ avec $\varphi(y) \in \overline{\{x'\}}$. Comme $x'$ est dans la fibre de $y'$, $x' \in \overline{\{\varphi(y')\}}$, et donc $\varphi(y) \in \overline{\{\varphi(y')\}}$.
\end{proof}

\begin{proof}[Démonstration de la proposition 1]
Les implications (i)~$\Leftrightarrow$~(ii), (ii)~$\Rightarrow$~(iii)~$\Leftrightarrow$~(iii bis) sont celles du lemme~\ref{lem:41-top}, (ii)~$\Rightarrow$~(iii) passant par (i). Il reste, pour $f$ un morphisme de schémas localement de présentation finie, à déduire (ii) de (iii bis). Or (iii bis) dit exactement que les générisations se relèvent le long de $f$, et un morphisme localement de présentation finie qui relève les générisations est ouvert : c'est un théorème classique (EGA~IV, 1.10.4), qui repose sur le théorème de Chevalley sur les parties constructibles. Les quatre conditions sont donc équivalentes.
\end{proof}

\begin{remarque}
L'implication (iii bis)~$\Rightarrow$~(ii) ne peut pas s'obtenir par la seule topologie. Soient $X = \mathbb N$ muni de la topologie discrète, $Y = \mathbb N$ muni de la topologie cofinie, et $f : X \to Y$ l'identité, qui est continue. Les fibres sont des points, et $\varphi = f^{-1}$ est une section générique. Les deux espaces étant séparés au sens T$_1$, $y \in \overline{\{y'\}}$ entraîne $y = y'$, et (iii bis) est triviale. Mais $f$ n'est pas ouverte : $f(\{0\}) = \{0\}$ n'est pas ouvert dans $Y$, puisque son complémentaire est infini. Ce n'est pas un morphisme de schémas ; l'exemple montre seulement que l'hypothèse de finitude, ou une autre, est indispensable à cette implication et à elle seule.
\end{remarque}

\begin{remarque}
Les implications (i)~$\Leftrightarrow$~(ii) et (i)~$\Rightarrow$~(iii)~$\Leftrightarrow$~(iii bis) valent pour toute application entre espaces topologiques munie d'une section générique, sans que $f$ soit de type fini, ni même continue.
\end{remarque}

\begin{remarque}
Seule l'implication (iii bis)~$\Rightarrow$~(ii) emploie l'hypothèse de finitude, sous la forme de la présentation finie locale, par le critère d'ouverture par relèvement des générisations (EGA~IV, 1.10.4), admis ici comme théorème classique.
\end{remarque}

\noindent Ne sont pas démontrés ici : le cas \q{par exemple $Y$ localement noethérien} (un morphisme de type fini sur une base localement noethérienne est localement de présentation finie) ; l'identification de la fibre schématique au sous-espace $f^{-1}(y)$ ; ce qui touche aux définitions~2 et~3 de la page~11 (morphisme universellement monogène, section monogène) ; la représentabilité du foncteur des sections monogènes (page~3), la réduction au cas d'une base locale complète (page~4), la construction des pages~5 et~6, le lemme de la page~7, le théorème des pages~8 et~9 et ses deux remarques, la question de la page~12 et le lemme barré de la page~13 ; la constructibilité, l'ouverture universelle et l'irréductibilité géométrique des fibres.
